\section{DSPy：声明式的 AI 编程}

\subsection{核心思想}

\begin{importantbox}{DSPy 的核心理念}
DSPy 让开发者用\textbf{Python 程序}定义 AI 系统的结构和逻辑，由框架自动\textbf{编译}出最优的 prompt 和 few-shot 示例。开发者只需声明``做什么''（what），不需要手工指定``怎么做''（how，即具体的 prompt 文本）。
\end{importantbox}

关键抽象：
\begin{itemize}
    \item \textbf{Signature}：定义模块的输入/输出类型（如 ``question -> answer''）
    \item \textbf{Module}：可组合的 AI 组件（如 ChainOfThought、ReAct）
    \item \textbf{Optimizer}：自动搜索最优 prompt、few-shot 示例和模块配置
    \item \textbf{Metric}：定义系统质量的评估标准
\end{itemize}

\subsection{优化器（Optimizers）}

DSPy 提供多种优化器来自动改进系统：

\begin{itemize}
    \item \textbf{BootstrapFewShot}：从训练数据中自动挑选最佳 few-shot 示例
    \item \textbf{COPRO}：自动优化 prompt 中的指令文本
    \item \textbf{MIPRO}：同时优化指令和示例
    \item \textbf{BootstrapFinetune}：生成训练数据并 fine-tune 底层模型
\end{itemize}

\subsection{编译 vs. 手写}

\begin{knowledgebox}{类比传统编程}
DSPy 之于 prompt engineering，如同编译器之于汇编语言。开发者用高级抽象编写程序，编译器（optimizer）自动生成底层实现（prompt）。这不仅提高了效率，还使系统具有可移植性——换模型只需重新编译。
\end{knowledgebox}

\subsection{本章小结}

DSPy 通过声明式编程和自动优化，将复合 AI 系统的构建从手工 prompt engineering 提升为系统化的软件工程。

